%% Section: Toward Precision Psychiatry — The Druggability Landscape
%% Position: between Therapeutic Landscape and Conclusions

\sectitle{Toward Precision Psychiatry: The Druggability Landscape}

The preceding sections established that emotional dysregulation spans
diagnostic boundaries, implicating a convergent set of molecular actors whose
dysfunction cuts across mood, anxiety, and trauma-related disorders. Yet
clinical practice has remained remarkably narrow in its pharmacological
response: the vast majority of patients presenting with transdiagnostic
emotional dysregulation receive a selective serotonin reuptake inhibitor as
first-line treatment. This section interrogates the full druggability landscape
of emotional dysregulation, drawing on systematic drug--gene interaction
databases to quantify the gap between what the genome offers and what the
clinic exploits. The picture that emerges is one of vast untapped therapeutic
space---space that precision psychiatry is uniquely positioned to claim.

\subsection*{Mapping the Druggable Genome of Emotional Dysregulation}

To characterise the pharmacological accessibility of the eight core genes
implicated in emotional dysregulation, we queried the Drug Gene Interaction
Database (DGIdb), a curated repository that aggregates drug--gene relationships
from over thirty source databases including PharmGKB, DrugBank, and the
Therapeutic Target Database \cite{freshour2021}. All eight genes---\gene{NR3C1},
\gene{BDNF}, \gene{SLC6A4}, \gene{COMT}, \gene{MAOA}, \gene{OXTR},
\gene{FKBP5}, and \gene{CRH}---reside within the DGIdb-defined druggable
genome, confirming that each encodes a protein whose structure and function
render it amenable to pharmacological modulation.

Collectively, these eight genes participate in 464 documented drug--gene
interactions, a number that reveals the sheer breadth of the actionable
pharmacological space surrounding emotional dysregulation. The distribution of
interactions across genes, however, is strikingly uneven and reflects both the
maturity of existing drug development programmes and the inherent tractability
of each protein class.

\gene{NR3C1}, encoding the glucocorticoid receptor, anchors the landscape with
134 drug--gene interactions---the largest share of any single gene. As a nuclear
hormone receptor and ligand-activated transcription factor, the glucocorticoid
receptor has long attracted pharmaceutical interest, predominantly through
glucocorticoid agonists developed for inflammatory and autoimmune conditions
\cite{oakley2013}. That the most interaction-rich gene in emotional
dysregulation belongs to the stress--hormone axis, rather than to the monoamine
systems that dominate current prescribing, underscores a fundamental mismatch
between pharmacological possibility and clinical habit.

\gene{BDNF}, encoding brain-derived neurotrophic factor, follows closely with
123 interactions. As a growth factor critical to synaptic plasticity and
neuronal survival, BDNF occupies a mechanistic position upstream of many
dysregulation phenotypes \cite{martinowich2007}. Several established
neuropsychiatric agents---ketamine, lithium, and valproic acid---exert their
therapeutic effects at least partly through BDNF-dependent signalling
\cite{duman2012}. The rapid antidepressant action of ketamine, mediated via
BDNF release in the prefrontal cortex, has reinvigorated interest in
neurotrophic mechanisms as direct therapeutic targets rather than incidental
downstream beneficiaries.

The serotonin transporter gene \gene{SLC6A4} accounts for 88 interactions,
encompassing the SSRIs and serotonin--norepinephrine reuptake inhibitors
(SNRIs) that constitute the current standard of care. While 88 interactions
represent a substantial pharmacological repertoire, this figure amounts to only
19\% of the total landscape. The remaining four genes---\gene{COMT} (42
interactions), \gene{MAOA} (39), \gene{OXTR} (18), \gene{FKBP5} (12), and
\gene{CRH} (8)---span catecholamine metabolism, oxytocinergic signalling, and
the hypothalamic--pituitary--adrenal (HPA) axis. Each harbours approved or
experimental agents: entacapone and tolcapone inhibit COMT-mediated catechol
degradation \cite{kiss2010}; phenelzine and tranylcypromine target monoamine
oxidase A; intranasal oxytocin and the antagonist atosiban modulate the
oxytocin receptor; SAFit2 selectively inhibits FKBP5 to restore glucocorticoid
receptor sensitivity \cite{gaali2015}; and antalarmin antagonises corticotropin-
releasing hormone receptor 1 to attenuate stress-axis hyperactivation
\cite{zorrilla2002}.

\subsection*{The Serotonin Bottleneck}

If the druggability landscape resembles a terrain, then current clinical
practice occupies a single well-trodden path through it. SSRIs and SNRIs
targeting \gene{SLC6A4} account for 88 of the 464 documented interactions, yet
these agents remain the default pharmacotherapy for emotional dysregulation
regardless of the patient's molecular profile. The remaining 376
interactions---spanning glucocorticoid, neurotrophic, catecholamine,
oxytocinergic, and CRH systems---constitute an untapped therapeutic landscape
that dwarfs the territory currently under clinical exploitation.

This serotonin bottleneck carries measurable clinical consequences. Large-scale
meta-analyses report that only 50--60\% of patients with major depression
achieve remission with first-line SSRI therapy, and response rates in
trauma-related and personality disorders are often lower still
\cite{cipriani2018}. Pharmacogenomic evidence increasingly explains this
variance. Stein and colleagues, in a landmark meta-analysis of \gene{SLC6A4}
genetic variation and antidepressant response, demonstrated that the
5-HTTLPR polymorphism in the promoter region of \gene{SLC6A4} significantly
modulates SSRI efficacy: carriers of the short allele exhibit reduced
serotonin transporter expression and, consequently, diminished response to
reuptake inhibition \cite{stein2021}. PharmGKB catalogues \gene{SLC6A4} as
PA312, mapping it to chromosome 17q11.2 and annotating multiple
pharmacogenomic associations that link transporter genotype to both efficacy
and adverse-effect profiles across SSRI classes.

These findings reframe treatment-resistant emotional dysregulation not as a
failure of pharmacology per se, but as a failure of targeting. The patient whose
dysregulation is driven primarily by HPA-axis hyperactivation via
\gene{NR3C1}--\gene{FKBP5} signalling, or by neurotrophic deficit mediated
through \gene{BDNF}, cannot reasonably be expected to respond to serotonin
reuptake blockade. Their molecular lesion lies elsewhere in the 464-interaction
landscape, and the appropriate pharmacological response lies there too.

\subsection*{Pathway-Level Evidence for Transdiagnostic Targeting}

The case for expanding beyond serotonin gains further support from curated
pathway analyses. WikiPathways catalogues \gene{SLC6A4} in ten curated
biological pathways, several of which illuminate the transdiagnostic
architecture of emotional dysregulation with unusual clarity. Pathway WP5420,
titled ``ADHD and autism (ASD) pathways,'' positions \gene{SLC6A4} within a
shared molecular circuit that spans attention-deficit, autistic, and affective
phenotypes---providing direct, curated evidence for the transdiagnostic
relevance of serotonergic signalling that epidemiological studies have long
suggested \cite{faraone2015}. That a single gene appears in a pathway
explicitly linking three diagnostic categories validates the broader thesis of
this review: emotional dysregulation is a dimensional construct whose molecular
substrates refuse to respect the boundaries drawn by categorical nosology.

Equally revealing, pathway WP1455 maps serotonin transporter activity per se,
while WP3594 situates \gene{SLC6A4} within circadian rhythm gene networks.
This latter connection is not merely academic. Sleep--wake dysregulation is among
the most consistent transdiagnostic features of emotional dysregulation, and
the placement of \gene{SLC6A4} at the intersection of serotonergic and
circadian circuitry suggests that chrono-pharmacological strategies---timing
drug administration to circadian phase, or targeting circadian regulators
directly---may offer therapeutic leverage that conventional fixed-dose SSRI
regimens cannot \cite{walker2020}.

These pathway-level data argue that precision psychiatry must move beyond
single-gene pharmacogenomics toward pathway-aware treatment selection. A
patient whose \gene{SLC6A4} genotype predicts poor SSRI response, and whose
clinical presentation features prominent HPA-axis markers (elevated cortisol
awakening response, stress-potentiated startle), may benefit more from agents
targeting the glucocorticoid receptor (\gene{NR3C1}) or its co-chaperone
\gene{FKBP5} than from dose escalation of an SSRI that their genome cannot
efficiently utilise.

\subsection*{A Framework for Pathway-Matched Pharmacotherapy}

The 464 drug--gene interactions mapped here suggest a framework in which
patients with emotional dysregulation are stratified not by DSM category but by
dominant molecular pathway. Under this framework,
five broad pharmacological axes emerge:

\begin{enumerate}
    \item \textbf{Serotonergic axis} (\gene{SLC6A4}, 88 interactions):
    SSRIs and SNRIs, appropriate when pharmacogenomic profiling confirms
    favourable transporter genotype and clinical features centre on
    rumination, anxiety sensitivity, and affective instability.

    \item \textbf{Glucocorticoid--stress axis} (\gene{NR3C1} + \gene{FKBP5} +
    \gene{CRH}, 154 interactions combined): Glucocorticoid receptor
    modulators, FKBP5 inhibitors such as SAFit2, and CRH antagonists such as
    antalarmin, appropriate when HPA-axis biomarkers indicate
    stress-system hyperactivation and trauma-related symptomatology dominates.

    \item \textbf{Neurotrophic axis} (\gene{BDNF}, 123 interactions):
    Ketamine and related glutamatergic agents, lithium, and valproic acid,
    appropriate when neurocognitive markers suggest plasticity deficits---for
    instance, impaired fear extinction learning or reduced hippocampal volume.

    \item \textbf{Catecholamine axis} (\gene{COMT} + \gene{MAOA}, 81
    interactions combined): COMT inhibitors and MAOIs, appropriate when
    executive-function deficits and dopaminergic or noradrenergic signatures
    predominate, as indexed by neuropsychological testing and catecholamine
    metabolite profiling.

    \item \textbf{Oxytocinergic axis} (\gene{OXTR}, 18 interactions):
    Intranasal oxytocin and related agents, appropriate when social
    cognitive deficits---impaired mentalising, reduced affiliative
    motivation---constitute the primary clinical concern.
\end{enumerate}

This five-axis model does not pretend to clinical readiness; considerable
translational work remains before pathway-matched prescribing can be validated
in randomised controlled trials. What the model does achieve is a
reorientation of the therapeutic imagination. Where current practice asks ``Has
the patient failed an SSRI?'' and proceeds to dose adjustment or augmentation
within the serotonergic axis, a precision framework asks ``Which of five
pharmacological axes best matches this patient's molecular and clinical
profile?'' The former question permits movement only along a single dimension;
the latter opens the full 464-interaction landscape to clinical navigation.

\subsection*{Barriers and the Path Forward}

Several obstacles stand between the present druggability map and its clinical
realisation. First, many of the 464 interactions involve agents developed for
non-psychiatric indications---glucocorticoid agonists for asthma, COMT
inhibitors for Parkinson's disease---whose repurposing for emotional
dysregulation will require dedicated efficacy and safety trials in psychiatric
populations. Second, the most pharmacologically novel targets (\gene{FKBP5},
\gene{CRH}) remain in early-phase development, with SAFit2 and antalarmin
lacking Phase~III data. Third, the clinical biomarkers required for
pathway-level stratification---cortisol dynamics, BDNF serum levels, 5-HTTLPR
genotyping---are not yet integrated into routine psychiatric assessment, and
their predictive validity for treatment selection requires prospective
validation \cite{bradley2020}.

Nevertheless, the convergence of three developments makes precision approaches
to emotional dysregulation increasingly tractable. Pharmacogenomic testing has
entered mainstream clinical practice, with combinatorial panels now covering
\gene{SLC6A4}, \gene{COMT}, and other dysregulation-relevant loci at declining
cost \cite{bousman2023}. High-throughput drug--gene interaction databases such
as DGIdb provide the systematic mapping needed to link genotype to actionable
pharmacology. And the transdiagnostic turn in psychiatric nosology---embodied
by the Research Domain Criteria (RDoC) framework and supported by the pathway
evidence reviewed here---creates the conceptual space in which pathway-matched
treatment can be scientifically evaluated and clinically implemented
\cite{insel2010}.

The druggability landscape of emotional dysregulation is not barren; it is
neglected. Eight genes yield 464 drug--gene interactions spanning five
pharmacological axes, yet clinical practice channels the vast majority of
patients through a single serotonergic corridor. The task ahead is not to
discover new drugs---though that too is welcome---but to match existing and
emerging agents to the patients most likely to benefit, guided by genetic,
epigenetic, and pathway-level evidence. In this sense, the path toward
precision psychiatry does not require a revolution in pharmacology. It requires
a revolution in targeting.
